\chapter{Amambay}\label{dept:Amambay}
\mapaparaguai{13}{Amambay}
\vfill\clearpage
\section{Amambay}

\begin{itemize}
\tightlist
\item
  Pedro Juan Caballero e a capital, polo de comercio de fronteira com
  Ponta Pora (Brasil); 106.721 hab.
\item
  Vocacao de pecuaria extensiva, soja e servicos de fronteira.
\item
  Lacuna de solar/pluviometria da capital fechada com dados NASA POWER
  (2001-2020).
\end{itemize}

\subsection{Ficha climática de Pedro Juan Caballero (capital
departamental)}

Fonte: NASA POWER Climatology API, período 2001-2020, coordenadas
-22.55°S / -55.73°W.

\textbf{Irradiância solar global --- (kWh/m²/dia):}

{\def\LTcaptype{none} % do not increment counter
\begin{longtable}[]{@{}
  >{\raggedright\arraybackslash}p{(\linewidth - 22\tabcolsep) * \real{0.0833}}
  >{\raggedright\arraybackslash}p{(\linewidth - 22\tabcolsep) * \real{0.0833}}
  >{\raggedright\arraybackslash}p{(\linewidth - 22\tabcolsep) * \real{0.0833}}
  >{\raggedright\arraybackslash}p{(\linewidth - 22\tabcolsep) * \real{0.0833}}
  >{\raggedright\arraybackslash}p{(\linewidth - 22\tabcolsep) * \real{0.0833}}
  >{\raggedright\arraybackslash}p{(\linewidth - 22\tabcolsep) * \real{0.0833}}
  >{\raggedright\arraybackslash}p{(\linewidth - 22\tabcolsep) * \real{0.0833}}
  >{\raggedright\arraybackslash}p{(\linewidth - 22\tabcolsep) * \real{0.0833}}
  >{\raggedright\arraybackslash}p{(\linewidth - 22\tabcolsep) * \real{0.0833}}
  >{\raggedright\arraybackslash}p{(\linewidth - 22\tabcolsep) * \real{0.0833}}
  >{\raggedright\arraybackslash}p{(\linewidth - 22\tabcolsep) * \real{0.0833}}
  >{\raggedright\arraybackslash}p{(\linewidth - 22\tabcolsep) * \real{0.0833}}@{}}
\toprule\noalign{}
\begin{minipage}[b]{\linewidth}\raggedright
Jan
\end{minipage} & \begin{minipage}[b]{\linewidth}\raggedright
Fev
\end{minipage} & \begin{minipage}[b]{\linewidth}\raggedright
Mar
\end{minipage} & \begin{minipage}[b]{\linewidth}\raggedright
Abr
\end{minipage} & \begin{minipage}[b]{\linewidth}\raggedright
Mai
\end{minipage} & \begin{minipage}[b]{\linewidth}\raggedright
Jun
\end{minipage} & \begin{minipage}[b]{\linewidth}\raggedright
Jul
\end{minipage} & \begin{minipage}[b]{\linewidth}\raggedright
Ago
\end{minipage} & \begin{minipage}[b]{\linewidth}\raggedright
Set
\end{minipage} & \begin{minipage}[b]{\linewidth}\raggedright
Out
\end{minipage} & \begin{minipage}[b]{\linewidth}\raggedright
Nov
\end{minipage} & \begin{minipage}[b]{\linewidth}\raggedright
Dez
\end{minipage} \\
\midrule\noalign{}
\endhead
\bottomrule\noalign{}
\endlastfoot
6.29 & 5.83 & 5.59 & 4.73 & 3.70 & 3.35 & 3.61 & 4.40 & 4.82 & 5.37 &
6.08 & 6.29 \\
\end{longtable}
}

Média anual: 5.00 kWh/m²/dia. Inclinação recomendada: 23° Norte (anual);
33° inverno; 13° verão.

\textbf{Precipitação --- (mm/mês):}

{\def\LTcaptype{none} % do not increment counter
\begin{longtable}[]{@{}
  >{\raggedright\arraybackslash}p{(\linewidth - 22\tabcolsep) * \real{0.0833}}
  >{\raggedright\arraybackslash}p{(\linewidth - 22\tabcolsep) * \real{0.0833}}
  >{\raggedright\arraybackslash}p{(\linewidth - 22\tabcolsep) * \real{0.0833}}
  >{\raggedright\arraybackslash}p{(\linewidth - 22\tabcolsep) * \real{0.0833}}
  >{\raggedright\arraybackslash}p{(\linewidth - 22\tabcolsep) * \real{0.0833}}
  >{\raggedright\arraybackslash}p{(\linewidth - 22\tabcolsep) * \real{0.0833}}
  >{\raggedright\arraybackslash}p{(\linewidth - 22\tabcolsep) * \real{0.0833}}
  >{\raggedright\arraybackslash}p{(\linewidth - 22\tabcolsep) * \real{0.0833}}
  >{\raggedright\arraybackslash}p{(\linewidth - 22\tabcolsep) * \real{0.0833}}
  >{\raggedright\arraybackslash}p{(\linewidth - 22\tabcolsep) * \real{0.0833}}
  >{\raggedright\arraybackslash}p{(\linewidth - 22\tabcolsep) * \real{0.0833}}
  >{\raggedright\arraybackslash}p{(\linewidth - 22\tabcolsep) * \real{0.0833}}@{}}
\toprule\noalign{}
\begin{minipage}[b]{\linewidth}\raggedright
Jan
\end{minipage} & \begin{minipage}[b]{\linewidth}\raggedright
Fev
\end{minipage} & \begin{minipage}[b]{\linewidth}\raggedright
Mar
\end{minipage} & \begin{minipage}[b]{\linewidth}\raggedright
Abr
\end{minipage} & \begin{minipage}[b]{\linewidth}\raggedright
Mai
\end{minipage} & \begin{minipage}[b]{\linewidth}\raggedright
Jun
\end{minipage} & \begin{minipage}[b]{\linewidth}\raggedright
Jul
\end{minipage} & \begin{minipage}[b]{\linewidth}\raggedright
Ago
\end{minipage} & \begin{minipage}[b]{\linewidth}\raggedright
Set
\end{minipage} & \begin{minipage}[b]{\linewidth}\raggedright
Out
\end{minipage} & \begin{minipage}[b]{\linewidth}\raggedright
Nov
\end{minipage} & \begin{minipage}[b]{\linewidth}\raggedright
Dez
\end{minipage} \\
\midrule\noalign{}
\endhead
\bottomrule\noalign{}
\endlastfoot
182 & 185 & 116 & 122 & 120 & 76 & 47 & 46 & 85 & 150 & 170 & 176 \\
\end{longtable}
}

Média anual: \textasciitilde1.476 mm/ano. Estação chuvosa Jan-Fev; seca
Jul-Ago.
\clearpage
\subsection*{Dados Departamentais}
\subsubsection{Desenvolvimento Humano e
Social}

\subsection{IDH Departamental}

{\def\LTcaptype{none} % do not increment counter
\begin{longtable}[]{@{}ll@{}}
\toprule\noalign{}
Indicador & Valor \\
\midrule\noalign{}
\endhead
\bottomrule\noalign{}
\endlastfoot
IDH & 0,720 \\
Classificação IDH & alto \\
Ranking nacional (1=melhor) & 5/18 \\
Esperança de vida & 73,6 anos \\
Anos de escolaridade média & 8,7 anos \\
RNB per capita (USD PPA) & 9.500 \\
\end{longtable}
}

\subsection{Indicadores de Pobreza}

{\def\LTcaptype{none} % do not increment counter
\begin{longtable}[]{@{}ll@{}}
\toprule\noalign{}
Indicador & Valor \\
\midrule\noalign{}
\endhead
\bottomrule\noalign{}
\endlastfoot
Taxa de pobreza total (\%) & 22,8\% \\
Taxa de pobreza extrema (\%) & 4,5\% \\
Índice de Gini & 0,448 \\
\end{longtable}
}

\subsection{Infraestrutura Básica}

{\def\LTcaptype{none} % do not increment counter
\begin{longtable}[]{@{}ll@{}}
\toprule\noalign{}
Indicador & Valor \\
\midrule\noalign{}
\endhead
\bottomrule\noalign{}
\endlastfoot
Acesso a água potável (\%) & 82,1\% \\
Acesso a saneamento (\%) & 80,5\% \\
\end{longtable}
}

\subsubsection{Segurança Pública}

\subsection{Indicadores}

{\def\LTcaptype{none} % do not increment counter
\begin{longtable}[]{@{}
  >{\raggedright\arraybackslash}p{(\linewidth - 2\tabcolsep) * \real{0.6111}}
  >{\raggedright\arraybackslash}p{(\linewidth - 2\tabcolsep) * \real{0.3889}}@{}}
\toprule\noalign{}
\begin{minipage}[b]{\linewidth}\raggedright
Indicador
\end{minipage} & \begin{minipage}[b]{\linewidth}\raggedright
Valor
\end{minipage} \\
\midrule\noalign{}
\endhead
\bottomrule\noalign{}
\endlastfoot
Taxa de homicídios (por 100k hab) & \textasciitilde30--50 (estimativa
2022--2024; histórico: 107,4 em 2016, 69,9 em 2017) \\
Crimes registrados (total/ano) & --- dados não desagregados
publicamente \\
Número de comisarías & \textasciitilde12 (estimativa departamental) \\
Índice segurança relativo & baixo \\
\end{longtable}
}

\subsubsection{Mercado de Terra Rural}

\subsection{Preços por Tipo de Terra}

{\def\LTcaptype{none} % do not increment counter
\begin{longtable}[]{@{}llll@{}}
\toprule\noalign{}
Tipo & Preço (USD/ha) & Faixa & Tendência \\
\midrule\noalign{}
\endhead
\bottomrule\noalign{}
\endlastfoot
Agrícola alta produtividade & 6.000 & 4.000 -- 9.000 & ↑ \\
Pastagem & 2.800 & 1.800 -- 4.000 & ↑ \\
Mata / reserva & 1.200 & 700 -- 2.000 & → \\
\end{longtable}
}

\subsection{Características do
Mercado}

{\def\LTcaptype{none} % do not increment counter
\begin{longtable}[]{@{}ll@{}}
\toprule\noalign{}
Parâmetro & Valor \\
\midrule\noalign{}
\endhead
\bottomrule\noalign{}
\endlastfoot
Valorização anual média (\%) & 7--9\% \\
Lote mínimo rural (ha) & 20 \\
Imposto territorial rural (\%) & 1\% sobre valor fiscal \\
Liquidez do mercado & média \\
\end{longtable}
}
\clearpage
\newpage
\secaoDiagnostico{Bella Vista}{\mapaDistrito{1302}}{Bella Vista Norte é um refúgio de baixo perfil situado em uma das
fronteiras mais emblemáticas do norte paraguaio. Oferece um isolamento
demográfico invejável e abundância de recursos naturais, funcionando
como um enclave de pecuária extensiva protegido pela geografia do Rio
Apa. Sua força reside no anonimato geográfico relativo e na capacidade
de subsistência baseada na proteína animal. Contudo, o contexto regional
de Amambay exige que o ocupante adote um perfil vigilante e autônomo,
mitigando os riscos de segurança tática da fronteira. É o local ideal
para projetos de autossuficiência que valorizem a privacidade total fora
dos eixos de fallout primários.

\subsubsection{Por que sim}

\begin{itemize}
\tightlist
\item
  Máximo isolamento demográfico em território vasto (privacidade).
\item
  Abundância de águas puras e solos aptos para pecuária e silvicultura.
\item
  Baixa visibilidade estratégica fora do contexto de fronteira.
\item
  Ambiente social pacífico nas comunidades rurais tradicionais.
\end{itemize}

\subsubsection{Por que não}

\begin{itemize}
\tightlist
\item
  Insegurança tática regional (narcotráfico e crime organizado no
  departamento).
\item
  Restrições da Lei de Fronteira (2532/05) para investidores
  estrangeiros.
\item
  Infraestrutura de saúde local limitada a atendimentos primários.
\end{itemize}

\subsubsection{Recomendação
Técnica}

Altamente indicado para estabelecimentos de autossuficiência pecuária
que busquem isolamento tático radical. Requer investimento em segurança
patrimonial, autonomia médica e hídrica. O GSS de 6.2 reflete o poder do
isolamento tático equilibrado pela vulnerabilidade institucional e
regional. Referencia atual de score: GSS 6.2.

\subsubsection{}

\begin{itemize}
\tightlist
\item
  \begin{enumerate}
  \def\labelenumi{\arabic{enumi})}
  \tightlist
  \item
    Dados censitários validados (INE 2022).
  \end{enumerate}
\item
  \begin{enumerate}
  \def\labelenumi{\arabic{enumi})}
  \setcounter{enumi}{1}
  \tightlist
  \item
    Relatórios de segurança pública e inteligência de fronteira
    (SENAD/MDI).
  \end{enumerate}
\item
  \begin{enumerate}
  \def\labelenumi{\arabic{enumi})}
  \setcounter{enumi}{2}
  \tightlist
  \item
    Mapa de solos e aptidão para invernada de Amambay (MAG).
  \end{enumerate}
\item
  \begin{enumerate}
  \def\labelenumi{\arabic{enumi})}
  \setcounter{enumi}{3}
  \tightlist
  \item
    Registro de pavimentação da Ruta PY08 (MOPC).
  \end{enumerate}
\end{itemize}}{dist:13_Amambay:Bella_Vista}
\begin{table}[ht!]
\small \begin{tabular}{p{3.0cm}p{0.8cm}p{6.0cm}} \toprule
\textbf{Critério} & \textbf{Nota} & \textbf{Justificativa} \\ \midrule
A - Ameaças Estratégicas & 5.5 & Ponto de fronteira internacional e hub histórico; visibilidade estratégica e riscos regionais de segurança \\
B - Risco Social & 7.5 & Densidade populacional extremamente baixa em território vasto; isolamento social satisfatório \\
C - Riscos Naturais & 7.0 & Geologia muito estável e baixo risco de grandes desastres naturais \\
D - Autossuficiência & 7.0 & Bons recursos de proteína animal e água; limitado pela distância de mercados técnicos \\
E - Institucional & 4.5 & Suporte institucional básico; dependência de centros regionais para serviços vitais \\
\midrule \textbf{GSS FINAL} & \textbf{6.2} & \textbf{Moderadamente Seguro (Ideal para quem busca isolamento rural produtivo com acesso fluvial, ciente dos desafios de segurança regional).} \\ \bottomrule \end{tabular}
\end{table}
\subsubsection{Vulnerabilidades e
Mitigação}\label{vuln-Bella_Vista-vulnerabilidades-e-mitigauxe7uxe3o}

\begin{itemize}
\tightlist
\item
  \textbf{Vulnerabilidade Tática de Fronteira:} Exposição a conflitos do
  crime organizado regional (Mitigação: residência rural afastada 10-15
  km da linha de fronteira seca/fluvial e investimento em segurança
  patrimonial robusta).
\item
  \textbf{Isolamento de Saúde:} Distância crítica de hospitais
  cirúrgicos (Mitigação: manutenção de estoque médico básico e
  treinamento em primeiros socorros avançados).
\item
  \textbf{Dependência de Rodovia Única:} Risco de bloqueio logístico
  (Mitigação: mapeamento de rotas rurais vicinais e posse de embarcação
  para via fluvial).
\end{itemize}

Notas (0-10): - A: 5.5 (Ponto de fronteira internacional e hub
histórico; visibilidade estratégica e riscos regionais de segurança) -
B: 7.5 (Densidade populacional extremamente baixa em território vasto;
isolamento social satisfatório) - C: 7.0 (Geologia muito estável e baixo
risco de grandes desastres naturais) - D: 7.0 (Bons recursos de proteína
animal e água; limitado pela distância de mercados técnicos) - E: 4.5
(Suporte institucional básico; dependência de centros regionais para
serviços vitais)

Calculo: - GSS = ((A * 2.5) + (B * 2.0) + (C * 1.5) + (D * 2.0) + (E *
2.0)) / 10 - GSS: 6.2

Classificacao: - Moderadamente Seguro (Ideal para quem busca isolamento
rural produtivo com acesso fluvial, ciente dos desafios de segurança
regional).

\begin{itemize}
\tightlist
\item
  Fonte MOPC infraestrutura: Reabilitação da Ruta PY08.
\end{itemize}
\subsubsection{Dossiê de Campo}\label{dossie-Bella_Vista-dossiuxea-de-campo}
\begin{description}[style=nextline,leftmargin=0.5cm,font=\bfseries]
\item[Coordenadas]  22°08'S 56°31'W.
\item[Topografia]  Relevo ondulado com extensas planícies de pastagem e áreas de mata nativa. Margeado ao norte pelo Rio Apa (fronteira natural com o Brasil). Altitude \textasciitilde 210m. Área vasta de \textasciitilde 3.901 km². Geologicamente muito estável, risco sísmico desprezível.
\item[Alvos Estratégicos]  Ponte Internacional sobre o Rio Apa (Conexão direta com Bela Vista, Brasil); Porto fluvial secundário; Locais históricos da Guerra da Tríplice Aliança nas proximidades. Ausência de bases militares massivas, mas zona de alta vigilância da Armada e Polícia de Fronteira devido à dinâmica transfronteiriça.
\item[Fallout]  Ventos predominantes do Norte (NE) e Leste. Baixíssimo risco de fallout direto de alvos nacionais primários.
\item[População]  12.027 habitantes (Censo 2022 Final). Perfil demográfico jovem, com economia integrada ao comércio de fronteira e pecuária.
\item[Densidade]  \textasciitilde 3,1 hab/km² (Densidade extremamente baixa, oferecendo excelente isolamento tático e privacidade fora do núcleo urbano consolidado).
\item[Vias de Acesso]  Conectividade asfáltica eficiente via ramal da Ruta PY08 (concluído em 2019). Localização que favorece a discrição por estar em um "final de linha" geográfico ao norte do departamento. Risco de isolamento tático se a ponte ou a rodovia PY08 forem bloqueadas.
\item[Serviços]  Centro de Saúde local e USF. Dependência de Pedro Juan Caballero (aprox. 100 km) para serviços hospitalares de alta complexidade. Infraestrutura comercial básica focada no intercâmbio fronteiriço.
\item[Custo de Vida]  Baixo; economia baseada na pecuária extensiva, erva-mate e comércio vicinal.
\item[Preço da Terra]  US\$ 2.500-5.000/ha (áreas desenvolvidas para pecuária); US\$ 1.500-2.500/ha (campos brutos/monte).
\item[Hidrologia]  Baixo risco de inundações sistêmicas devido à topografia ondulada. O Rio Apa possui praias fluviais e regime de cheias previsível. Presença de diversos arroios tributários de alta qualidade hídrica.
\item[Clima]  Tropical Úmido. Verões intensos e chuvosos; invernos amenos.
\item[Energia]  Rede nacional (Itaipu); instável em áreas rurais remotas durante temporais.
\item[Água]  Abundante via Rio Apa, nascentes naturais e poços artesianos produtivos. Excelente disponibilidade hídrica subterrânea.
\item[Qualidade do Solo]  Solo franco-arenoso fértil; aptidão para pastagens naturais, silvicultura e cultivo de erva-mate. Áreas mecanizadas em expansão.
\item[Recursos Locais]  Grande produção de gado bovino e erva-mate. Elevadíssimo potencial para autossuficiência de proteína animal em nível de propriedade devido à vasta área produtiva e baixa população.
\item[Segurança]  Zona de fronteira com vigilância tática permanente. Criminalidade urbana comum moderada, mas sensível a atividades ilícitas transfronteiriças (narcotráfico/contrabando) típicas do departamento de Amambay. Coesão social baseada na cultura agropecuária tradicional.
\item[Leis Local: Município consolidado com forte intercâmbio cultural com o Brasil. Restrição de Fronteira (Lei 2532/05)]  Aplicável a terras rurais na faixa de 50km da fronteira brasileira.
\end{description}
\textbf{Fonte climática:} NASA POWER Climatology API (período 2001-2020)

\textbf{Fonte luminosa:} estimativa world atlas

\paragraph{Irradiação Solar
(kWh/m²/dia)}\label{irradiauxe7uxe3o-solar-kwhmuxb2dia}

{\def\LTcaptype{none} % do not increment counter
\begin{longtable}[]{@{}
  >{\raggedright\arraybackslash}p{(\linewidth - 24\tabcolsep) * \real{0.0746}}
  >{\raggedright\arraybackslash}p{(\linewidth - 24\tabcolsep) * \real{0.0746}}
  >{\raggedright\arraybackslash}p{(\linewidth - 24\tabcolsep) * \real{0.0746}}
  >{\raggedright\arraybackslash}p{(\linewidth - 24\tabcolsep) * \real{0.0746}}
  >{\raggedright\arraybackslash}p{(\linewidth - 24\tabcolsep) * \real{0.0746}}
  >{\raggedright\arraybackslash}p{(\linewidth - 24\tabcolsep) * \real{0.0746}}
  >{\raggedright\arraybackslash}p{(\linewidth - 24\tabcolsep) * \real{0.0746}}
  >{\raggedright\arraybackslash}p{(\linewidth - 24\tabcolsep) * \real{0.0746}}
  >{\raggedright\arraybackslash}p{(\linewidth - 24\tabcolsep) * \real{0.0746}}
  >{\raggedright\arraybackslash}p{(\linewidth - 24\tabcolsep) * \real{0.0746}}
  >{\raggedright\arraybackslash}p{(\linewidth - 24\tabcolsep) * \real{0.0746}}
  >{\raggedright\arraybackslash}p{(\linewidth - 24\tabcolsep) * \real{0.0746}}
  >{\raggedright\arraybackslash}p{(\linewidth - 24\tabcolsep) * \real{0.1045}}@{}}
\toprule\noalign{}
\begin{minipage}[b]{\linewidth}\raggedright
Jan
\end{minipage} & \begin{minipage}[b]{\linewidth}\raggedright
Fev
\end{minipage} & \begin{minipage}[b]{\linewidth}\raggedright
Mar
\end{minipage} & \begin{minipage}[b]{\linewidth}\raggedright
Abr
\end{minipage} & \begin{minipage}[b]{\linewidth}\raggedright
Mai
\end{minipage} & \begin{minipage}[b]{\linewidth}\raggedright
Jun
\end{minipage} & \begin{minipage}[b]{\linewidth}\raggedright
Jul
\end{minipage} & \begin{minipage}[b]{\linewidth}\raggedright
Ago
\end{minipage} & \begin{minipage}[b]{\linewidth}\raggedright
Set
\end{minipage} & \begin{minipage}[b]{\linewidth}\raggedright
Out
\end{minipage} & \begin{minipage}[b]{\linewidth}\raggedright
Nov
\end{minipage} & \begin{minipage}[b]{\linewidth}\raggedright
Dez
\end{minipage} & \begin{minipage}[b]{\linewidth}\raggedright
Média
\end{minipage} \\
\midrule\noalign{}
\endhead
\bottomrule\noalign{}
\endlastfoot
6.44 & 5.92 & 5.66 & 4.74 & 3.62 & 3.27 & 3.51 & 4.32 & 4.78 & 5.46 &
6.22 & 6.37 & \textbf{5.03} \\
\end{longtable}
}

\textbf{Inclinação solar recomendada:} 22° N (anual) · 32° N (inverno
jun-ago) · 12° N (verão nov-jan)

\paragraph{Precipitação (mm/mês)}\label{precipitauxe7uxe3o-mmmuxeas}

{\def\LTcaptype{none} % do not increment counter
\begin{longtable}[]{@{}
  >{\raggedright\arraybackslash}p{(\linewidth - 24\tabcolsep) * \real{0.0704}}
  >{\raggedright\arraybackslash}p{(\linewidth - 24\tabcolsep) * \real{0.0704}}
  >{\raggedright\arraybackslash}p{(\linewidth - 24\tabcolsep) * \real{0.0704}}
  >{\raggedright\arraybackslash}p{(\linewidth - 24\tabcolsep) * \real{0.0704}}
  >{\raggedright\arraybackslash}p{(\linewidth - 24\tabcolsep) * \real{0.0704}}
  >{\raggedright\arraybackslash}p{(\linewidth - 24\tabcolsep) * \real{0.0704}}
  >{\raggedright\arraybackslash}p{(\linewidth - 24\tabcolsep) * \real{0.0704}}
  >{\raggedright\arraybackslash}p{(\linewidth - 24\tabcolsep) * \real{0.0704}}
  >{\raggedright\arraybackslash}p{(\linewidth - 24\tabcolsep) * \real{0.0704}}
  >{\raggedright\arraybackslash}p{(\linewidth - 24\tabcolsep) * \real{0.0704}}
  >{\raggedright\arraybackslash}p{(\linewidth - 24\tabcolsep) * \real{0.0704}}
  >{\raggedright\arraybackslash}p{(\linewidth - 24\tabcolsep) * \real{0.0704}}
  >{\raggedright\arraybackslash}p{(\linewidth - 24\tabcolsep) * \real{0.1549}}@{}}
\toprule\noalign{}
\begin{minipage}[b]{\linewidth}\raggedright
Jan
\end{minipage} & \begin{minipage}[b]{\linewidth}\raggedright
Fev
\end{minipage} & \begin{minipage}[b]{\linewidth}\raggedright
Mar
\end{minipage} & \begin{minipage}[b]{\linewidth}\raggedright
Abr
\end{minipage} & \begin{minipage}[b]{\linewidth}\raggedright
Mai
\end{minipage} & \begin{minipage}[b]{\linewidth}\raggedright
Jun
\end{minipage} & \begin{minipage}[b]{\linewidth}\raggedright
Jul
\end{minipage} & \begin{minipage}[b]{\linewidth}\raggedright
Ago
\end{minipage} & \begin{minipage}[b]{\linewidth}\raggedright
Set
\end{minipage} & \begin{minipage}[b]{\linewidth}\raggedright
Out
\end{minipage} & \begin{minipage}[b]{\linewidth}\raggedright
Nov
\end{minipage} & \begin{minipage}[b]{\linewidth}\raggedright
Dez
\end{minipage} & \begin{minipage}[b]{\linewidth}\raggedright
Total/ano
\end{minipage} \\
\midrule\noalign{}
\endhead
\bottomrule\noalign{}
\endlastfoot
166 & 161 & 105 & 108 & 109 & 64 & 40 & 36 & 74 & 133 & 159 & 154 &
\textbf{1315 mm} \\
\end{longtable}
}

\paragraph{Poluição Luminosa}\label{poluiuxe7uxe3o-luminosa}

{\def\LTcaptype{none} % do not increment counter
\begin{longtable}[]{@{}ll@{}}
\toprule\noalign{}
Parâmetro & Valor \\
\midrule\noalign{}
\endhead
\bottomrule\noalign{}
\endlastfoot
Escala Bortle & 4 --- Céu rural-suburbano \\
Radiância artificial & 2.0 nW/cm²/sr \\
\end{longtable}
}
\paragraph{Solo (SoilGrids 2.0, média ponderada 0--30
cm)}

{\def\LTcaptype{none} % do not increment counter
\begin{longtable}[]{@{}ll@{}}
\toprule\noalign{}
Parâmetro & Valor \\
\midrule\noalign{}
\endhead
\bottomrule\noalign{}
\endlastfoot
pH (H\(_2\)O) & 5.72 \\
Carbono orgânico (SOC) & 13.0 g/kg \\
Argila & 15.48 \% \\
Areia & 67.22 \% \\
Silte (calc.) & 17.3 \% \\
Densidade aparente & 1.32 g/cm³ \\
\textbf{Aptidão agrícola} & \textbf{Média} \\
\end{longtable}
}

Fonte: ISRIC SoilGrids 2.0 via WCS. Coords: -22.1333°, -56.5167°. Média
ponderada camadas 0-5, 5-15, 15-30 cm.
\subsubsection{Indicadores Sociais}

\textbf{Fontes:} PNUD 2020, INE Censo 2022, INE EPHC 2023, Ministerio
Público 2024\\
\textbf{Nota:} Valores marcados como \emph{(dept.)} referem-se ao
departamento de Amambay; valores \emph{(dist.)} são específicos deste
distrito. Dados marcados \emph{(est.)} são estimativas.

{\def\LTcaptype{none} % do not increment counter
\begin{longtable}[]{@{}
  >{\raggedright\arraybackslash}p{(\linewidth - 4\tabcolsep) * \real{0.4231}}
  >{\raggedright\arraybackslash}p{(\linewidth - 4\tabcolsep) * \real{0.2692}}
  >{\raggedright\arraybackslash}p{(\linewidth - 4\tabcolsep) * \real{0.3077}}@{}}
\toprule\noalign{}
\begin{minipage}[b]{\linewidth}\raggedright
Indicador
\end{minipage} & \begin{minipage}[b]{\linewidth}\raggedright
Valor
\end{minipage} & \begin{minipage}[b]{\linewidth}\raggedright
Âmbito
\end{minipage} \\
\midrule\noalign{}
\endhead
\bottomrule\noalign{}
\endlastfoot
IDH (2020) & 0,720 & dept. \\
Ranking IDH nacional & 5/18 & dept. \\
Esperança de vida & 73,6 anos & dept. \\
Escolaridade média & 7.0 anos & dist. \\
RNB per capita (USD PPA) & 9.500 & dept. \\
Pobreza monetária (\%) & 22,8\% & dept. \\
Pobreza extrema (\%) & 4,5\% & dept. \\
Índice de Gini & 0,448 & dept. \\
Acesso a água potável (\%) & 82,1\% & dept. \\
Acesso a saneamento (\%) & 80,5\% & dept. \\
Taxa de homicídios (est., /100k hab) & \textasciitilde30--50 (estimativa
2022--2024; histórico: 107,4 em 2016, 69,9 em 2017) & dept. \\
Índice de segurança & baixo & dept. \\
População (Censo 2022) & 10.447 hab. & dist. \\
Idade mediana & 29.0 anos & dist. \\
Taxa de fecundidade & N/D & dist. \\
\end{longtable}
}
\subsubsection{Combustível}

\textbf{Referência:} PETROPAR / postos locais (2024) \textbf{Tipo de
localidade:} Interior

{\def\LTcaptype{none} % do not increment counter
\begin{longtable}[]{@{}lll@{}}
\toprule\noalign{}
Combustível & USD/litro & Gs/litro (aprox.) \\
\midrule\noalign{}
\endhead
\bottomrule\noalign{}
\endlastfoot
Gasolina 93 oct & 1.00 & 7,400 \\
Gasolina 97 oct (premium) & 1.10 & 8,140 \\
Diesel & 0.93 & 6,882 \\
\end{longtable}
}

\begin{quote}
Preços podem variar ±5\% conforme posto e sazonalidade. Chaco e interior
remoto apresentam maior variação.
\end{quote}

\subsubsection{Cobertura Celular}

\textbf{Fonte:} CONATEL PY / operadoras (2024)

{\def\LTcaptype{none} % do not increment counter
\begin{longtable}[]{@{}ll@{}}
\toprule\noalign{}
Parâmetro & Valor \\
\midrule\noalign{}
\endhead
\bottomrule\noalign{}
\endlastfoot
Cobertura 4G população (dept.) & 85\% \\
Cobertura 4G área rural & 68\% \\
Melhor operadora & Tigo \\
Qualidade rural & moderada \\
\end{longtable}
}

\begin{quote}
Para áreas rurais fora do núcleo urbano, recomenda-se chip Tigo como
principal e Personal como backup.
\end{quote}

\subsubsection{Internet}

\textbf{Fonte:} CONATEL / Speedtest Ookla (2024)

{\def\LTcaptype{none} % do not increment counter
\begin{longtable}[]{@{}ll@{}}
\toprule\noalign{}
Parâmetro & Valor \\
\midrule\noalign{}
\endhead
\bottomrule\noalign{}
\endlastfoot
Velocidade média download & 55 Mbps \\
Domicílios com internet (dept.) & 62\% \\
Tecnologia predominante & rádio/fibra \\
Opção rural & Starlink disponível (\textasciitilde USD 44/mês) \\
\end{longtable}
}

\subsubsection{Mercado Imobiliário e Terra
Rural}

\textbf{Fonte:} INDERT / Clasificados.com.py (2024)

{\def\LTcaptype{none} % do not increment counter
\begin{longtable}[]{@{}ll@{}}
\toprule\noalign{}
Tipo & Referência \\
\midrule\noalign{}
\endhead
\bottomrule\noalign{}
\endlastfoot
Terra agrícola alta prod. (USD/ha) & 5,000 \\
Imóvel urbano (USD/m²) & 900 \\
Aluguel 2 quartos (USD/mês) & 380 \\
\end{longtable}
}

\begin{quote}
Valores de referência departamental. Localidades menores podem ter
preços 20--40\% abaixo da capital departamental.
\end{quote}

\subsubsection{Saúde}

\textbf{Fonte:} MSPBS / IPS Paraguay (2024-2026), consolidação
departamental e proxy local

{\def\LTcaptype{none} % do not increment counter
\begin{longtable}[]{@{}ll@{}}
\toprule\noalign{}
Serviço & Disponibilidade \\
\midrule\noalign{}
\endhead
\bottomrule\noalign{}
\endlastfoot
USF / Posto de Saúde & sim \\
Hospital Regional & sim \\
IPS (seguro social) & não \\
Farmácia & sim \\
Distância ao hospital de referência & \textasciitilde120 km (Pedro Juan
Caballero) \\
\end{longtable}
}

\textbf{Principais estabelecimentos:} USF local; referencia hospitalar
em Pedro Juan Caballero

\textbf{Observação para imigrantes:} Atendimento primario local ou em
raio curto; casos de maior complexidade seguem para Pedro Juan
Caballero. Cobertura privada continua recomendada para especialidades e
urgências de maior complexidade.
\newpage
\secaoDiagnostico{Capitan Bado}{\mapaDistrito{1303}}{Capitán Bado é um polo de recursos de elite situado em uma das zonas
mais taticamente complexas do Paraguai. Oferece solos de produtividade
mundial e um isolamento demográfico que garante privacidade absoluta.
Sua força reside na excepcional abundância de alimentos e águas.
Contudo, a proximidade com a fronteira seca e o contexto regional de
Amambay exigem que o ocupante adote um perfil vigilante e autônomo. É o
local ideal para quem busca terra de alta performance e anonimato
geográfico, desde que possua meios próprios de segurança e suporte
logístico.

\subsubsection{Por que sim}

\begin{itemize}
\tightlist
\item
  Solos de altíssima produtividade e clima favorável.
\item
  Baixíssima densidade populacional (privacidade absoluta).
\item
  Abundância de recursos hídricos e proteicos.
\item
  Longe de eixos de fallout tático metropolitano.
\end{itemize}

\subsubsection{Por que não}

\begin{itemize}
\tightlist
\item
  Insegurança tática regional latente (fronteira).
\item
  Restrições da Lei de Fronteira para investidores estrangeiros.
\item
  Distância de centros urbanos com serviços de alta complexidade.
\end{itemize}

\subsubsection{Recomendação
Técnica}

Altamente indicado para estabelecimentos de autossuficiência tecnificada
que valorizem a qualidade do solo. Requer investimento em segurança
patrimonial e autonomia em saúde. O GSS de 7.1 reflete a superioridade
dos recursos produtivos em equilíbrio com os desafios do contexto
regional. Referencia atual de score: GSS 7.1.

\subsubsection{}

\begin{itemize}
\tightlist
\item
  \begin{enumerate}
  \def\labelenumi{\arabic{enumi})}
  \tightlist
  \item
    Dados censitários validados (INE 2022).
  \end{enumerate}
\item
  \begin{enumerate}
  \def\labelenumi{\arabic{enumi})}
  \setcounter{enumi}{1}
  \tightlist
  \item
    Relatórios de produção agrícola departamental (MAG).
  \end{enumerate}
\item
  \begin{enumerate}
  \def\labelenumi{\arabic{enumi})}
  \setcounter{enumi}{2}
  \tightlist
  \item
    Mapa de solos e bacias hidrográficas de Amambay (MAG).
  \end{enumerate}
\item
  \begin{enumerate}
  \def\labelenumi{\arabic{enumi})}
  \setcounter{enumi}{3}
  \tightlist
  \item
    Registro de segurança pública e inteligência de fronteira.
  \end{enumerate}
\end{itemize}}{dist:13_Amambay:Capitan_Bado}
\begin{table}[ht!]
\small \begin{tabular}{p{3.0cm}p{0.8cm}p{6.0cm}} \toprule
\textbf{Critério} & \textbf{Nota} & \textbf{Justificativa} \\ \midrule
A - Ameaças Estratégicas & 6.5 & Localização remota; excelente invisibilidade tática fora da fronteira seca, equilibrada por riscos regionais \\
B - Risco Social & 7.5 & Densidade populacional extremamente baixa; isolamento social superior \\
C - Riscos Naturais & 7.0 & Geologia muito estável e relevo serrano seguro contra inundações \\
D - Autossuficiência & 8.0 & Recursos agrícolas de elite e abundância hídrica natural \\
E - Institucional & 6.5 & Estabilidade institucional presente, mas sob vigilância tática de fronteira \\
\midrule \textbf{GSS FINAL} & \textbf{7.1} & \textbf{Seguro (Altamente recomendado para quem busca autossuficiência em solo de elite e isolamento tático, ciente dos desafios de segurança regional).} \\ \bottomrule \end{tabular}
\end{table}
\subsubsection{Vulnerabilidades e
Mitigação}\label{vuln-Capitan_Bado-vulnerabilidades-e-mitigauxe7uxe3o}

\begin{itemize}
\tightlist
\item
  \textbf{Vulnerabilidade Tática Regional:} Insegurança ligada ao
  narcotráfico regional (Mitigação: investimento em segurança
  patrimonial robusta e residência rural afastada da linha de
  fronteira).
\item
  \textbf{Isolamento de Saúde:} Distância significativa de hospitais
  especializados (Mitigação: manutenção de estoque médico básico e kit
  trauma local).
\item
  \textbf{Dependência de Redes:} Fragilidade elétrica em tempestades
  (Mitigação: instalação de sistemas solares fotovoltaicos potentes).
\end{itemize}

Notas (0-10): - A: 6.5 (Localização remota; excelente invisibilidade
tática fora da fronteira seca, equilibrada por riscos regionais) - B:
7.5 (Densidade populacional extremamente baixa; isolamento social
superior) - C: 7.0 (Geologia muito estável e relevo serrano seguro
contra inundações) - D: 8.0 (Recursos agrícolas de elite e abundância
hídrica natural) - E: 6.5 (Estabilidade institucional presente, mas sob
vigilância tática de fronteira)

Calculo: - GSS = ((A * 2.5) + (B * 2.0) + (C * 1.5) + (D * 2.0) + (E *
2.0)) / 10 - GSS: 7.1

Classificacao: - Seguro (Altamente recomendado para quem busca
autossuficiência em solo de elite e isolamento tático, ciente dos
desafios de segurança regional).
\subsubsection{Dossiê de Campo}\label{dossie-Capitan_Bado-dossiuxea-de-campo}
\begin{description}[style=nextline,leftmargin=0.5cm,font=\bfseries]
\item[Coordenadas]  23°15'S 55°32'W.
\item[Topografia]  Relevo ondulado com presença de serras (Cordillera de Amambay) e solos de elite (Terra Roxa). Altitude \textasciitilde 480m. Área vasta de \textasciitilde 3.205 km². Geologicamente estável.
\item[Alvos Estratégicos]  Polo de produção agrícola e ervateira; Ponto de fronteira seca internacional com Coronel Sapucaia (Brasil). Localização estratégica para o controle do nordeste paraguaio, mas situada em zona de alta complexidade tática regional.
\item[Fallout]  Ventos predominantes do Norte (NE) e Leste. Baixo risco de fallout direto de alvos continentais primários devido à altitude e isolamento.
\item[População]  18.851 habitantes (Censo 2022 Final). População com forte componente demográfico rural e indígena, integrada ao agronegócio e comércio fronteiriço.
\item[Densidade]  \textasciitilde 5,9 hab/km² (Densidade extremamente baixa, oferecendo isolamento tático e privacidade fora do núcleo urbano).
\item[Vias de Acesso]  Conectividade asfáltica via Ruta PY11. Localização em "final de linha" geográfico que favorece o anonimato estratégico, mas dificulta o apoio logístico rápido da capital nacional.
\item[Serviços]  Centro de Saúde local e IPS. Dependência de Pedro Juan Caballero (\textasciitilde 110 km) para atendimentos de alta complexidade. Infraestrutura comercial focada em suprimentos agrícolas.
\item[Custo de Vida]  Baixo; economia baseada na pecuária, grãos e comércio de fronteira.
\item[Preço da Terra]  US\$ 6.000-10.000/ha (terras mecanizadas); US\$ 2.000-4.000/ha (áreas de pecuária).
\item[Hidrologia]  Baixo risco de inundações sistêmicas devido ao relevo serrano e drenagem natural eficiente. Diversas nascentes e arroios de alta qualidade hídrica.
\item[Clima]  Subtropical Úmido. Verões quentes e chuvosos; invernos frescos com geadas frequentes que favorecem o controle de pragas agrícolas.
\item[Energia]  Rede nacional (Itaipu); instável em áreas rurais remotas.
\item[Água]  Abundância hídrica via poços artesianos e nascentes na serra; excelente qualidade subterrânea.
\item[Qualidade do Solo]  Latossolo Vermelho de Elite; solo profundo e extremamente fértil. Excelente aptidão para soja, milho e erva-mate.
\item[Recursos Locais]  Grande produção de grãos e gado bovino. Elevadíssimo potencial para autossuficiência alimentar básica em nível de propriedade devido à vasta área produtiva e baixa população.
\item[Segurança]  Zona de fronteira com vigilância tática elevada. Historicamente sensível a atividades do crime organizado transfronteiriço. Coesão social baseada na cultura produtora tradicional. Segurança institucional presente, mas desafiada pela extensão territorial.
\item[Leis Local: Município consolidado. Restrição de Fronteira (Lei 2532/05)]  Aplicável a terras rurais na faixa de 50km da fronteira brasileira.
\end{description}
\textbf{Fonte climática:} NASA POWER Climatology API (período 2001-2020)

\textbf{Fonte luminosa:} estimativa world atlas

\paragraph{Irradiação Solar
(kWh/m²/dia)}\label{irradiauxe7uxe3o-solar-kwhmuxb2dia}

{\def\LTcaptype{none} % do not increment counter
\begin{longtable}[]{@{}
  >{\raggedright\arraybackslash}p{(\linewidth - 24\tabcolsep) * \real{0.0746}}
  >{\raggedright\arraybackslash}p{(\linewidth - 24\tabcolsep) * \real{0.0746}}
  >{\raggedright\arraybackslash}p{(\linewidth - 24\tabcolsep) * \real{0.0746}}
  >{\raggedright\arraybackslash}p{(\linewidth - 24\tabcolsep) * \real{0.0746}}
  >{\raggedright\arraybackslash}p{(\linewidth - 24\tabcolsep) * \real{0.0746}}
  >{\raggedright\arraybackslash}p{(\linewidth - 24\tabcolsep) * \real{0.0746}}
  >{\raggedright\arraybackslash}p{(\linewidth - 24\tabcolsep) * \real{0.0746}}
  >{\raggedright\arraybackslash}p{(\linewidth - 24\tabcolsep) * \real{0.0746}}
  >{\raggedright\arraybackslash}p{(\linewidth - 24\tabcolsep) * \real{0.0746}}
  >{\raggedright\arraybackslash}p{(\linewidth - 24\tabcolsep) * \real{0.0746}}
  >{\raggedright\arraybackslash}p{(\linewidth - 24\tabcolsep) * \real{0.0746}}
  >{\raggedright\arraybackslash}p{(\linewidth - 24\tabcolsep) * \real{0.0746}}
  >{\raggedright\arraybackslash}p{(\linewidth - 24\tabcolsep) * \real{0.1045}}@{}}
\toprule\noalign{}
\begin{minipage}[b]{\linewidth}\raggedright
Jan
\end{minipage} & \begin{minipage}[b]{\linewidth}\raggedright
Fev
\end{minipage} & \begin{minipage}[b]{\linewidth}\raggedright
Mar
\end{minipage} & \begin{minipage}[b]{\linewidth}\raggedright
Abr
\end{minipage} & \begin{minipage}[b]{\linewidth}\raggedright
Mai
\end{minipage} & \begin{minipage}[b]{\linewidth}\raggedright
Jun
\end{minipage} & \begin{minipage}[b]{\linewidth}\raggedright
Jul
\end{minipage} & \begin{minipage}[b]{\linewidth}\raggedright
Ago
\end{minipage} & \begin{minipage}[b]{\linewidth}\raggedright
Set
\end{minipage} & \begin{minipage}[b]{\linewidth}\raggedright
Out
\end{minipage} & \begin{minipage}[b]{\linewidth}\raggedright
Nov
\end{minipage} & \begin{minipage}[b]{\linewidth}\raggedright
Dez
\end{minipage} & \begin{minipage}[b]{\linewidth}\raggedright
Média
\end{minipage} \\
\midrule\noalign{}
\endhead
\bottomrule\noalign{}
\endlastfoot
6.76 & 6.20 & 5.42 & 4.44 & 3.36 & 2.83 & 3.21 & 3.93 & 4.61 & 5.49 &
6.42 & 6.78 & \textbf{4.95} \\
\end{longtable}
}

\textbf{Inclinação solar recomendada:} 25° N (anual) · 35° N (inverno
jun-ago) · 15° N (verão nov-jan)

\paragraph{Precipitação (mm/mês)}\label{precipitauxe7uxe3o-mmmuxeas}

{\def\LTcaptype{none} % do not increment counter
\begin{longtable}[]{@{}
  >{\raggedright\arraybackslash}p{(\linewidth - 24\tabcolsep) * \real{0.0704}}
  >{\raggedright\arraybackslash}p{(\linewidth - 24\tabcolsep) * \real{0.0704}}
  >{\raggedright\arraybackslash}p{(\linewidth - 24\tabcolsep) * \real{0.0704}}
  >{\raggedright\arraybackslash}p{(\linewidth - 24\tabcolsep) * \real{0.0704}}
  >{\raggedright\arraybackslash}p{(\linewidth - 24\tabcolsep) * \real{0.0704}}
  >{\raggedright\arraybackslash}p{(\linewidth - 24\tabcolsep) * \real{0.0704}}
  >{\raggedright\arraybackslash}p{(\linewidth - 24\tabcolsep) * \real{0.0704}}
  >{\raggedright\arraybackslash}p{(\linewidth - 24\tabcolsep) * \real{0.0704}}
  >{\raggedright\arraybackslash}p{(\linewidth - 24\tabcolsep) * \real{0.0704}}
  >{\raggedright\arraybackslash}p{(\linewidth - 24\tabcolsep) * \real{0.0704}}
  >{\raggedright\arraybackslash}p{(\linewidth - 24\tabcolsep) * \real{0.0704}}
  >{\raggedright\arraybackslash}p{(\linewidth - 24\tabcolsep) * \real{0.0704}}
  >{\raggedright\arraybackslash}p{(\linewidth - 24\tabcolsep) * \real{0.1549}}@{}}
\toprule\noalign{}
\begin{minipage}[b]{\linewidth}\raggedright
Jan
\end{minipage} & \begin{minipage}[b]{\linewidth}\raggedright
Fev
\end{minipage} & \begin{minipage}[b]{\linewidth}\raggedright
Mar
\end{minipage} & \begin{minipage}[b]{\linewidth}\raggedright
Abr
\end{minipage} & \begin{minipage}[b]{\linewidth}\raggedright
Mai
\end{minipage} & \begin{minipage}[b]{\linewidth}\raggedright
Jun
\end{minipage} & \begin{minipage}[b]{\linewidth}\raggedright
Jul
\end{minipage} & \begin{minipage}[b]{\linewidth}\raggedright
Ago
\end{minipage} & \begin{minipage}[b]{\linewidth}\raggedright
Set
\end{minipage} & \begin{minipage}[b]{\linewidth}\raggedright
Out
\end{minipage} & \begin{minipage}[b]{\linewidth}\raggedright
Nov
\end{minipage} & \begin{minipage}[b]{\linewidth}\raggedright
Dez
\end{minipage} & \begin{minipage}[b]{\linewidth}\raggedright
Total/ano
\end{minipage} \\
\midrule\noalign{}
\endhead
\bottomrule\noalign{}
\endlastfoot
126 & 139 & 135 & 149 & 143 & 74 & 63 & 38 & 73 & 163 & 190 & 179 &
\textbf{1477 mm} \\
\end{longtable}
}

\paragraph{Poluição Luminosa}\label{poluiuxe7uxe3o-luminosa}

{\def\LTcaptype{none} % do not increment counter
\begin{longtable}[]{@{}ll@{}}
\toprule\noalign{}
Parâmetro & Valor \\
\midrule\noalign{}
\endhead
\bottomrule\noalign{}
\endlastfoot
Escala Bortle & 4 --- Céu rural-suburbano \\
Radiância artificial & 2.0 nW/cm²/sr \\
\end{longtable}
}
\paragraph{Solo (SoilGrids 2.0, média ponderada 0--30
cm)}

{\def\LTcaptype{none} % do not increment counter
\begin{longtable}[]{@{}ll@{}}
\toprule\noalign{}
Parâmetro & Valor \\
\midrule\noalign{}
\endhead
\bottomrule\noalign{}
\endlastfoot
pH (H\(_2\)O) & 5.1 \\
Carbono orgânico (SOC) & 15.42 g/kg \\
Argila & 24.02 \% \\
Areia & 63.89 \% \\
Silte (calc.) & 12.1 \% \\
Densidade aparente & 1.25 g/cm³ \\
\textbf{Aptidão agrícola} & \textbf{Média} \\
\end{longtable}
}

Fonte: ISRIC SoilGrids 2.0 via WCS. Coords: -23.25°, -55.5333°. Média
ponderada camadas 0-5, 5-15, 15-30 cm.

\begin{itemize}
\tightlist
\item
  \textbf{Homicidios/100k:} \textasciitilde25,0 (Regional/Fronteira).
\end{itemize}
\subsubsection{Indicadores Sociais}

\textbf{Fontes:} PNUD 2020, INE Censo 2022, INE EPHC 2023, Ministerio
Público 2024\\
\textbf{Nota:} Valores marcados como \emph{(dept.)} referem-se ao
departamento de Amambay; valores \emph{(dist.)} são específicos deste
distrito. Dados marcados \emph{(est.)} são estimativas.

{\def\LTcaptype{none} % do not increment counter
\begin{longtable}[]{@{}
  >{\raggedright\arraybackslash}p{(\linewidth - 4\tabcolsep) * \real{0.4231}}
  >{\raggedright\arraybackslash}p{(\linewidth - 4\tabcolsep) * \real{0.2692}}
  >{\raggedright\arraybackslash}p{(\linewidth - 4\tabcolsep) * \real{0.3077}}@{}}
\toprule\noalign{}
\begin{minipage}[b]{\linewidth}\raggedright
Indicador
\end{minipage} & \begin{minipage}[b]{\linewidth}\raggedright
Valor
\end{minipage} & \begin{minipage}[b]{\linewidth}\raggedright
Âmbito
\end{minipage} \\
\midrule\noalign{}
\endhead
\bottomrule\noalign{}
\endlastfoot
IDH (2020) & 0,720 & dept. \\
Ranking IDH nacional & 5/18 & dept. \\
Esperança de vida & 73,6 anos & dept. \\
Escolaridade média & 6.9 anos & dist. \\
RNB per capita (USD PPA) & 9.500 & dept. \\
Pobreza monetária (\%) & 22,8\% & dept. \\
Pobreza extrema (\%) & 4,5\% & dept. \\
Índice de Gini & 0,448 & dept. \\
Acesso a água potável (\%) & 82,1\% & dept. \\
Acesso a saneamento (\%) & 80,5\% & dept. \\
Taxa de homicídios (est., /100k hab) & \textasciitilde30--50 (estimativa
2022--2024; histórico: 107,4 em 2016, 69,9 em 2017) & dept. \\
Índice de segurança & baixo & dept. \\
População (Censo 2022) & 18.851 habitantes & dist. \\
Idade mediana & N/D & dist. \\
Taxa de fecundidade & N/D & dist. \\
\end{longtable}
}
\subsubsection{Combustível}

\textbf{Referência:} PETROPAR / postos locais (2024) \textbf{Tipo de
localidade:} Interior

{\def\LTcaptype{none} % do not increment counter
\begin{longtable}[]{@{}lll@{}}
\toprule\noalign{}
Combustível & USD/litro & Gs/litro (aprox.) \\
\midrule\noalign{}
\endhead
\bottomrule\noalign{}
\endlastfoot
Gasolina 93 oct & 1.00 & 7,400 \\
Gasolina 97 oct (premium) & 1.10 & 8,140 \\
Diesel & 0.93 & 6,882 \\
\end{longtable}
}

\begin{quote}
Preços podem variar ±5\% conforme posto e sazonalidade. Chaco e interior
remoto apresentam maior variação.
\end{quote}

\subsubsection{Cobertura Celular}

\textbf{Fonte:} CONATEL PY / operadoras (2024)

{\def\LTcaptype{none} % do not increment counter
\begin{longtable}[]{@{}ll@{}}
\toprule\noalign{}
Parâmetro & Valor \\
\midrule\noalign{}
\endhead
\bottomrule\noalign{}
\endlastfoot
Cobertura 4G população (dept.) & 85\% \\
Cobertura 4G área rural & 68\% \\
Melhor operadora & Tigo \\
Qualidade rural & moderada \\
\end{longtable}
}

\begin{quote}
Para áreas rurais fora do núcleo urbano, recomenda-se chip Tigo como
principal e Personal como backup.
\end{quote}

\subsubsection{Internet}

\textbf{Fonte:} CONATEL / Speedtest Ookla (2024)

{\def\LTcaptype{none} % do not increment counter
\begin{longtable}[]{@{}ll@{}}
\toprule\noalign{}
Parâmetro & Valor \\
\midrule\noalign{}
\endhead
\bottomrule\noalign{}
\endlastfoot
Velocidade média download & 55 Mbps \\
Domicílios com internet (dept.) & 62\% \\
Tecnologia predominante & rádio/fibra \\
Opção rural & Starlink disponível (\textasciitilde USD 44/mês) \\
\end{longtable}
}

\subsubsection{Mercado Imobiliário e Terra
Rural}

\textbf{Fonte:} INDERT / Clasificados.com.py (2024)

{\def\LTcaptype{none} % do not increment counter
\begin{longtable}[]{@{}ll@{}}
\toprule\noalign{}
Tipo & Referência \\
\midrule\noalign{}
\endhead
\bottomrule\noalign{}
\endlastfoot
Terra agrícola alta prod. (USD/ha) & 5,000 \\
Imóvel urbano (USD/m²) & 900 \\
Aluguel 2 quartos (USD/mês) & 380 \\
\end{longtable}
}

\begin{quote}
Valores de referência departamental. Localidades menores podem ter
preços 20--40\% abaixo da capital departamental.
\end{quote}

\subsubsection{Saúde}

\textbf{Fonte:} MSPBS / IPS Paraguay (2024-2026), consolidação
departamental e proxy local

{\def\LTcaptype{none} % do not increment counter
\begin{longtable}[]{@{}ll@{}}
\toprule\noalign{}
Serviço & Disponibilidade \\
\midrule\noalign{}
\endhead
\bottomrule\noalign{}
\endlastfoot
USF / Posto de Saúde & sim \\
Hospital Regional & não \\
IPS (seguro social) & sim \\
Farmácia & sim \\
Distância ao hospital de referência & \textasciitilde455 km (Pedro Juan
Caballero) \\
\end{longtable}
}

\textbf{Principais estabelecimentos:} USF local; referencia hospitalar
em Pedro Juan Caballero

\textbf{Observação para imigrantes:} Atendimento primario local ou em
raio curto; casos de maior complexidade seguem para Pedro Juan
Caballero. Cobertura privada continua recomendada para especialidades e
urgências de maior complexidade.
\newpage
\secaoDiagnostico{Cerro Cora}{\mapaDistrito{1306}}{Cerro Corá é um dos refúgios táticos mais interessantes do nordeste
paraguaio. Oferece um isolamento geográfico e demográfico de alto nível,
protegido pelo relevo serrano e pela vasta área do Parque Nacional. Sua
força absoluta reside na abundância de recursos hídricos e na segurança
social sustentada por sua importância histórica. É o local ideal para
quem busca privacidade total em solo fértil, sem abdicar da
conectividade rodoviária de elite (PY05). O contexto regional exige
autonomia e vigilância, transformando a geografia acidentada em uma
barreira natural contra instabilidades metropolitanas.

\subsubsection{Por que sim}

\begin{itemize}
\tightlist
\item
  Máximo isolamento geográfico em território vasto e serrano.
\item
  Abundância de águas puras e recursos proteicos (gado).
\item
  Localização estratégica ``fora do radar'' militar primário.
\item
  Sem restrições da Lei de Fronteira.
\end{itemize}

\subsubsection{Por que não}

\begin{itemize}
\tightlist
\item
  Insegurança tática regional latente no departamento de Amambay.
\item
  Dependência de Pedro Juan Caballero para serviços de saúde complexos.
\item
  Pequena escala comercial local exige planejamento logístico de
  suprimentos.
\end{itemize}

\subsubsection{Recomendação
Técnica}

Altamente indicado para estabelecimentos de autossuficiência familiar de
alto padrão que valorizem o isolamento natural. Requer investimento em
autonomia energética e segurança patrimonial. O GSS de 7.6 reflete a
excepcional solidez do isolamento geográfico reforçada pela base de
recursos naturais. Referencia atual de score: GSS 7.6.

\subsubsection{}

\begin{itemize}
\tightlist
\item
  \begin{enumerate}
  \def\labelenumi{\arabic{enumi})}
  \tightlist
  \item
    Dados censitários validados (INE 2022).
  \end{enumerate}
\item
  \begin{enumerate}
  \def\labelenumi{\arabic{enumi})}
  \setcounter{enumi}{1}
  \tightlist
  \item
    Relatórios de monitoramento ambiental do Parque Nacional Cerro Corá
    (MADES).
  \end{enumerate}
\item
  \begin{enumerate}
  \def\labelenumi{\arabic{enumi})}
  \setcounter{enumi}{2}
  \tightlist
  \item
    Mapa de solos e hidrografia de Amambay (MAG).
  \end{enumerate}
\item
  \begin{enumerate}
  \def\labelenumi{\arabic{enumi})}
  \setcounter{enumi}{3}
  \tightlist
  \item
    Registro de segurança pública departamental.
  \end{enumerate}
\end{itemize}}{dist:13_Amambay:Cerro_Cora}
\begin{table}[ht!]
\small \begin{tabular}{p{3.0cm}p{0.8cm}p{6.0cm}} \toprule
\textbf{Critério} & \textbf{Nota} & \textbf{Justificativa} \\ \midrule
A - Ameaças Estratégicas & 7.5 & Isolamento tático estratégico; relevo serrano defensivo e ausência de alvos militares globais \\
B - Risco Social & 8.0 & Densidade populacional extremamente baixa; isolamento social de alto nível \\
C - Riscos Naturais & 7.5 & Geologia muito estável e relevo seguro contra inundações sistêmicas \\
D - Autossuficiência & 8.0 & Excelentes recursos hídricos e de proteína animal; suporte ambiental do parque \\
E - Institucional & 7.0 & Forte segurança institucional histórica e suporte em expansão como novo município \\
\midrule \textbf{GSS FINAL} & \textbf{7.6} & \textbf{Seguro (Localidade de elite para quem busca isolamento tático em ambiente histórico e natural, ciente dos desafios regionais).} \\ \bottomrule \end{tabular}
\end{table}
\subsubsection{Vulnerabilidades e
Mitigação}\label{vuln-Cerro_Cora-vulnerabilidades-e-mitigauxe7uxe3o}

\begin{itemize}
\tightlist
\item
  \textbf{Vulnerabilidade Tática Regional:} Insegurança ligada ao
  narcotráfico no departamento (Mitigação: residência rural protegida
  pelo relevo e investimento em segurança patrimonial).
\item
  \textbf{Isolamento de Saúde:} Necessidade de deslocamento para PJC em
  emergências (Mitigação: manutenção de estoque médico básico e kit
  trauma local).
\item
  \textbf{Pressão Turística:} Visibilidade tática em feriados nacionais
  devido ao Parque Cerro Corá (Mitigação: residência rural afastada 5-10
  km do perímetro do parque histórico).
\end{itemize}

Notas (0-10): - A: 7.5 (Isolamento tático estratégico; relevo serrano
defensivo e ausência de alvos militares globais) - B: 8.0 (Densidade
populacional extremamente baixa; isolamento social de alto nível) - C:
7.5 (Geologia muito estável e relevo seguro contra inundações
sistêmicas) - D: 8.0 (Excelentes recursos hídricos e de proteína animal;
suporte ambiental do parque) - E: 7.0 (Forte segurança institucional
histórica e suporte em expansão como novo município)

Calculo: - GSS = ((A * 2.5) + (B * 2.0) + (C * 1.5) + (D * 2.0) + (E *
2.0)) / 10 - GSS: 7.6

Classificacao: - Seguro (Localidade de elite para quem busca isolamento
tático em ambiente histórico e natural, ciente dos desafios regionais).

\begin{itemize}
\tightlist
\item
  Fonte MOPC infraestrutura: Reabilitação da Ruta PY05.
\end{itemize}
\subsubsection{Dossiê de Campo}\label{dossie-Cerro_Cora-dossiuxea-de-campo}
\begin{description}[style=nextline,leftmargin=0.5cm,font=\bfseries]
\item[Coordenadas]  22°39'S 56°01'W.
\item[Topografia]  Relevo ondulado a acidentado, caracterizado por colinas e serras (Cordillera de Amambay) e vales férteis. Altitude \textasciitilde 350m. Área massiva de \textasciitilde 1.442 km². Geologicamente estável.
\item[Alvos Estratégicos]  Parque Nacional Cerro Corá (Local de imenso valor histórico e simbólico nacional; reserva ambiental estratégica); Ponto de conexão na Ruta PY05 (Eixo Concepción-PJC). Ausência de bases militares massivas, mas zona de alta visibilidade histórica e vigilância tática regional.
\item[Fallout]  Ventos predominantes do Norte (NE) e Leste. Baixíssimo risco de fallout direto; relevo serrano oferece barreiras naturais parciais.
\item[População]  10.690 habitantes (Censo 2022 Final). Município de criação recente (2020), consolidando áreas rurais e comunidades originárias.
\item[Densidade]  \textasciitilde 7,4 hab/km² (Densidade baixíssima, oferecendo excelente isolamento tático e privacidade absoluta fora da vila urbana central).
\item[Vias de Acesso]  Localização privilegiada sobre a Ruta PY05 (pavimentada). Conectividade de elite para Pedro Juan Caballero (45 km) e Concepción. Facilidade logística tática em tempos de paz, com rotas vicinais que permitem o desaparecimento rural.
\item[Serviços]  Unidade de Saúde da Família (USF) local. Dependência direta de Pedro Juan Caballero para serviços hospitalares de alta complexidade e centros comerciais diversificados.
\item[Custo de Vida]  Muito baixo; economia baseada na pecuária, pequenas plantações e ecoturismo histórico.
\item[Preço da Terra]  US\$ 3.000-6.000/ha (áreas desenvolvidas); US\$ 1.500-3.000/ha (áreas de monte/pecuária).
\item[Hidrologia]  Baixo risco de inundações sistêmicas devido ao relevo favorável à drenagem. Abundância de arroios perenes (ex: Rio Aquidabán) e nascentes de alta qualidade hídrica.
\item[Clima]  Subtropical Úmido. Verões intensos e chuvosos; invernos frescos com geadas ocasionais. Microclima influenciado pela vegetação do parque nacional.
\item[Energia]  Rede nacional (Itaipu); instável em áreas rurais remotas.
\item[Água]  Abundante via nascentes naturais e poços artesianos; excelente qualidade hídrica superficial e subterrânea.
\item[Qualidade do Solo]  Solo franco-arenoso a argiloso fértil; aptidão para pastagens naturais, silvicultura e agricultura de subsistência.
\item[Recursos Locais]  Grande produção de gado bovino e recursos florestais. Elevadíssimo potencial para autossuficiência calórica familiar baseada em proteína animal e horticultura protegida.
\item[Segurança]  Zona rural pacífica em termos de crime comum. No entanto, exige atenção tática regional devido ao contexto de Amambay (narcotráfico e trânsito de ilícitos nas proximidades de PJC). Forte identidade histórica nacionalista. Estabilidade social sustentada pela cultura tradicional.
\item[Leis Local: Município de criação recente (ex-Chiriguelo). Livre de Restrição de Fronteira]  Localizado fora da faixa de 50km da fronteira seca internacional (protegido pela linha de PJC).
\end{description}
\textbf{Fonte climática:} NASA POWER Climatology API (período 2001-2020)

\textbf{Fonte luminosa:} estimativa world atlas

\paragraph{Irradiação Solar
(kWh/m²/dia)}\label{irradiauxe7uxe3o-solar-kwhmuxb2dia}

{\def\LTcaptype{none} % do not increment counter
\begin{longtable}[]{@{}
  >{\raggedright\arraybackslash}p{(\linewidth - 24\tabcolsep) * \real{0.0746}}
  >{\raggedright\arraybackslash}p{(\linewidth - 24\tabcolsep) * \real{0.0746}}
  >{\raggedright\arraybackslash}p{(\linewidth - 24\tabcolsep) * \real{0.0746}}
  >{\raggedright\arraybackslash}p{(\linewidth - 24\tabcolsep) * \real{0.0746}}
  >{\raggedright\arraybackslash}p{(\linewidth - 24\tabcolsep) * \real{0.0746}}
  >{\raggedright\arraybackslash}p{(\linewidth - 24\tabcolsep) * \real{0.0746}}
  >{\raggedright\arraybackslash}p{(\linewidth - 24\tabcolsep) * \real{0.0746}}
  >{\raggedright\arraybackslash}p{(\linewidth - 24\tabcolsep) * \real{0.0746}}
  >{\raggedright\arraybackslash}p{(\linewidth - 24\tabcolsep) * \real{0.0746}}
  >{\raggedright\arraybackslash}p{(\linewidth - 24\tabcolsep) * \real{0.0746}}
  >{\raggedright\arraybackslash}p{(\linewidth - 24\tabcolsep) * \real{0.0746}}
  >{\raggedright\arraybackslash}p{(\linewidth - 24\tabcolsep) * \real{0.0746}}
  >{\raggedright\arraybackslash}p{(\linewidth - 24\tabcolsep) * \real{0.1045}}@{}}
\toprule\noalign{}
\begin{minipage}[b]{\linewidth}\raggedright
Jan
\end{minipage} & \begin{minipage}[b]{\linewidth}\raggedright
Fev
\end{minipage} & \begin{minipage}[b]{\linewidth}\raggedright
Mar
\end{minipage} & \begin{minipage}[b]{\linewidth}\raggedright
Abr
\end{minipage} & \begin{minipage}[b]{\linewidth}\raggedright
Mai
\end{minipage} & \begin{minipage}[b]{\linewidth}\raggedright
Jun
\end{minipage} & \begin{minipage}[b]{\linewidth}\raggedright
Jul
\end{minipage} & \begin{minipage}[b]{\linewidth}\raggedright
Ago
\end{minipage} & \begin{minipage}[b]{\linewidth}\raggedright
Set
\end{minipage} & \begin{minipage}[b]{\linewidth}\raggedright
Out
\end{minipage} & \begin{minipage}[b]{\linewidth}\raggedright
Nov
\end{minipage} & \begin{minipage}[b]{\linewidth}\raggedright
Dez
\end{minipage} & \begin{minipage}[b]{\linewidth}\raggedright
Média
\end{minipage} \\
\midrule\noalign{}
\endhead
\bottomrule\noalign{}
\endlastfoot
6.41 & 5.92 & 5.60 & 4.65 & 3.53 & 3.14 & 3.44 & 4.20 & 4.70 & 5.36 &
6.22 & 6.44 & \textbf{4.97} \\
\end{longtable}
}

\textbf{Inclinação solar recomendada:} 23° N (anual) · 33° N (inverno
jun-ago) · 13° N (verão nov-jan)

\paragraph{Precipitação (mm/mês)}\label{precipitauxe7uxe3o-mmmuxeas}

{\def\LTcaptype{none} % do not increment counter
\begin{longtable}[]{@{}
  >{\raggedright\arraybackslash}p{(\linewidth - 24\tabcolsep) * \real{0.0704}}
  >{\raggedright\arraybackslash}p{(\linewidth - 24\tabcolsep) * \real{0.0704}}
  >{\raggedright\arraybackslash}p{(\linewidth - 24\tabcolsep) * \real{0.0704}}
  >{\raggedright\arraybackslash}p{(\linewidth - 24\tabcolsep) * \real{0.0704}}
  >{\raggedright\arraybackslash}p{(\linewidth - 24\tabcolsep) * \real{0.0704}}
  >{\raggedright\arraybackslash}p{(\linewidth - 24\tabcolsep) * \real{0.0704}}
  >{\raggedright\arraybackslash}p{(\linewidth - 24\tabcolsep) * \real{0.0704}}
  >{\raggedright\arraybackslash}p{(\linewidth - 24\tabcolsep) * \real{0.0704}}
  >{\raggedright\arraybackslash}p{(\linewidth - 24\tabcolsep) * \real{0.0704}}
  >{\raggedright\arraybackslash}p{(\linewidth - 24\tabcolsep) * \real{0.0704}}
  >{\raggedright\arraybackslash}p{(\linewidth - 24\tabcolsep) * \real{0.0704}}
  >{\raggedright\arraybackslash}p{(\linewidth - 24\tabcolsep) * \real{0.0704}}
  >{\raggedright\arraybackslash}p{(\linewidth - 24\tabcolsep) * \real{0.1549}}@{}}
\toprule\noalign{}
\begin{minipage}[b]{\linewidth}\raggedright
Jan
\end{minipage} & \begin{minipage}[b]{\linewidth}\raggedright
Fev
\end{minipage} & \begin{minipage}[b]{\linewidth}\raggedright
Mar
\end{minipage} & \begin{minipage}[b]{\linewidth}\raggedright
Abr
\end{minipage} & \begin{minipage}[b]{\linewidth}\raggedright
Mai
\end{minipage} & \begin{minipage}[b]{\linewidth}\raggedright
Jun
\end{minipage} & \begin{minipage}[b]{\linewidth}\raggedright
Jul
\end{minipage} & \begin{minipage}[b]{\linewidth}\raggedright
Ago
\end{minipage} & \begin{minipage}[b]{\linewidth}\raggedright
Set
\end{minipage} & \begin{minipage}[b]{\linewidth}\raggedright
Out
\end{minipage} & \begin{minipage}[b]{\linewidth}\raggedright
Nov
\end{minipage} & \begin{minipage}[b]{\linewidth}\raggedright
Dez
\end{minipage} & \begin{minipage}[b]{\linewidth}\raggedright
Total/ano
\end{minipage} \\
\midrule\noalign{}
\endhead
\bottomrule\noalign{}
\endlastfoot
159 & 176 & 118 & 134 & 132 & 82 & 55 & 49 & 92 & 171 & 188 & 179 &
\textbf{1542 mm} \\
\end{longtable}
}

\paragraph{Poluição Luminosa}\label{poluiuxe7uxe3o-luminosa}

{\def\LTcaptype{none} % do not increment counter
\begin{longtable}[]{@{}ll@{}}
\toprule\noalign{}
Parâmetro & Valor \\
\midrule\noalign{}
\endhead
\bottomrule\noalign{}
\endlastfoot
Escala Bortle & 4 --- Céu rural-suburbano \\
Radiância artificial & 2.0 nW/cm²/sr \\
\end{longtable}
}
\paragraph{Solo (SoilGrids 2.0, média ponderada 0--30
cm)}

{\def\LTcaptype{none} % do not increment counter
\begin{longtable}[]{@{}ll@{}}
\toprule\noalign{}
Parâmetro & Valor \\
\midrule\noalign{}
\endhead
\bottomrule\noalign{}
\endlastfoot
pH (H\(_2\)O) & 5.4 \\
Carbono orgânico (SOC) & 19.14 g/kg \\
Argila & 26.11 \% \\
Areia & 50.87 \% \\
Silte (calc.) & 23.0 \% \\
Densidade aparente & 1.24 g/cm³ \\
\textbf{Aptidão agrícola} & \textbf{Média} \\
\end{longtable}
}

Fonte: ISRIC SoilGrids 2.0 via WCS. Coords: -22.65°, -56.0167°. Média
ponderada camadas 0-5, 5-15, 15-30 cm.
\subsubsection{Indicadores Sociais}

\textbf{Fontes:} PNUD 2020, INE Censo 2022, INE EPHC 2023, Ministerio
Público 2024\\
\textbf{Nota:} Valores marcados como \emph{(dept.)} referem-se ao
departamento de Amambay; valores \emph{(dist.)} são específicos deste
distrito. Dados marcados \emph{(est.)} são estimativas.

{\def\LTcaptype{none} % do not increment counter
\begin{longtable}[]{@{}
  >{\raggedright\arraybackslash}p{(\linewidth - 4\tabcolsep) * \real{0.4231}}
  >{\raggedright\arraybackslash}p{(\linewidth - 4\tabcolsep) * \real{0.2692}}
  >{\raggedright\arraybackslash}p{(\linewidth - 4\tabcolsep) * \real{0.3077}}@{}}
\toprule\noalign{}
\begin{minipage}[b]{\linewidth}\raggedright
Indicador
\end{minipage} & \begin{minipage}[b]{\linewidth}\raggedright
Valor
\end{minipage} & \begin{minipage}[b]{\linewidth}\raggedright
Âmbito
\end{minipage} \\
\midrule\noalign{}
\endhead
\bottomrule\noalign{}
\endlastfoot
IDH (2020) & 0,720 & dept. \\
Ranking IDH nacional & 5/18 & dept. \\
Esperança de vida & 73,6 anos & dept. \\
Escolaridade média & 6.2 anos & dist. \\
RNB per capita (USD PPA) & 9.500 & dept. \\
Pobreza monetária (\%) & 22,8\% & dept. \\
Pobreza extrema (\%) & 4,5\% & dept. \\
Índice de Gini & 0,448 & dept. \\
Acesso a água potável (\%) & 82,1\% & dept. \\
Acesso a saneamento (\%) & 80,5\% & dept. \\
Taxa de homicídios (est., /100k hab) & \textasciitilde30--50 (estimativa
2022--2024; histórico: 107,4 em 2016, 69,9 em 2017) & dept. \\
Índice de segurança & baixo & dept. \\
População (Censo 2022) & 10.690 habitantes & dist. \\
Idade mediana & N/D & dist. \\
Taxa de fecundidade & N/D & dist. \\
\end{longtable}
}
\subsubsection{Combustível}

\textbf{Referência:} PETROPAR / postos locais (2024) \textbf{Tipo de
localidade:} Interior

{\def\LTcaptype{none} % do not increment counter
\begin{longtable}[]{@{}lll@{}}
\toprule\noalign{}
Combustível & USD/litro & Gs/litro (aprox.) \\
\midrule\noalign{}
\endhead
\bottomrule\noalign{}
\endlastfoot
Gasolina 93 oct & 1.00 & 7,400 \\
Gasolina 97 oct (premium) & 1.10 & 8,140 \\
Diesel & 0.93 & 6,882 \\
\end{longtable}
}

\begin{quote}
Preços podem variar ±5\% conforme posto e sazonalidade. Chaco e interior
remoto apresentam maior variação.
\end{quote}

\subsubsection{Cobertura Celular}

\textbf{Fonte:} CONATEL PY / operadoras (2024)

{\def\LTcaptype{none} % do not increment counter
\begin{longtable}[]{@{}ll@{}}
\toprule\noalign{}
Parâmetro & Valor \\
\midrule\noalign{}
\endhead
\bottomrule\noalign{}
\endlastfoot
Cobertura 4G população (dept.) & 85\% \\
Cobertura 4G área rural & 68\% \\
Melhor operadora & Tigo \\
Qualidade rural & moderada \\
\end{longtable}
}

\begin{quote}
Para áreas rurais fora do núcleo urbano, recomenda-se chip Tigo como
principal e Personal como backup.
\end{quote}

\subsubsection{Internet}

\textbf{Fonte:} CONATEL / Speedtest Ookla (2024)

{\def\LTcaptype{none} % do not increment counter
\begin{longtable}[]{@{}ll@{}}
\toprule\noalign{}
Parâmetro & Valor \\
\midrule\noalign{}
\endhead
\bottomrule\noalign{}
\endlastfoot
Velocidade média download & 55 Mbps \\
Domicílios com internet (dept.) & 62\% \\
Tecnologia predominante & rádio/fibra \\
Opção rural & Starlink disponível (\textasciitilde USD 44/mês) \\
\end{longtable}
}

\subsubsection{Mercado Imobiliário e Terra
Rural}

\textbf{Fonte:} INDERT / Clasificados.com.py (2024)

{\def\LTcaptype{none} % do not increment counter
\begin{longtable}[]{@{}ll@{}}
\toprule\noalign{}
Tipo & Referência \\
\midrule\noalign{}
\endhead
\bottomrule\noalign{}
\endlastfoot
Terra agrícola alta prod. (USD/ha) & 5,000 \\
Imóvel urbano (USD/m²) & 900 \\
Aluguel 2 quartos (USD/mês) & 380 \\
\end{longtable}
}

\begin{quote}
Valores de referência departamental. Localidades menores podem ter
preços 20--40\% abaixo da capital departamental.
\end{quote}

\subsubsection{Saúde}

\textbf{Fonte:} MSPBS / IPS Paraguay (2024-2026), consolidação
departamental e proxy local

{\def\LTcaptype{none} % do not increment counter
\begin{longtable}[]{@{}ll@{}}
\toprule\noalign{}
Serviço & Disponibilidade \\
\midrule\noalign{}
\endhead
\bottomrule\noalign{}
\endlastfoot
USF / Posto de Saúde & sim \\
Hospital Regional & sim \\
IPS (seguro social) & não \\
Farmácia & sim \\
Distância ao hospital de referência & \textasciitilde100 km (Pedro Juan
Caballero) \\
\end{longtable}
}

\textbf{Principais estabelecimentos:} USF local; referencia hospitalar
em Pedro Juan Caballero

\textbf{Observação para imigrantes:} Atendimento primario local ou em
raio curto; casos de maior complexidade seguem para Pedro Juan
Caballero. Cobertura privada continua recomendada para especialidades e
urgências de maior complexidade.
\newpage
\secaoDiagnostico{Karapai}{\mapaDistrito{1305}}{Karapaí é um polo de alta produtividade agrícola em um dos departamentos
mais sensíveis do Paraguai. Oferece solos de elite mundial e um
isolamento demográfico que garante privacidade quase absoluta. Sua força
reside na riqueza de recursos naturais e na nova conectividade
asfáltica. Contudo, o contexto de segurança do departamento de Amambay
exige que o ocupante adote um perfil vigilante e autônomo. É o local
ideal para quem busca terra de alta performance fora dos grandes eixos
de fallout, desde que mitigados os riscos de isolamento de serviços e as
dinâmicas de segurança regional.

\subsubsection{Por que sim}

\begin{itemize}
\tightlist
\item
  Solos de elite (Terra Roxa) e abundância de águas puras.
\item
  Baixíssima densidade populacional (privacidade absoluta).
\item
  Conectividade asfáltica moderna via Ruta PY11.
\item
  Fora da zona de restrição de fronteira.
\end{itemize}

\subsubsection{Por que não}

\begin{itemize}
\tightlist
\item
  Insegurança tática regional ligada ao narcotráfico no departamento.
\item
  Infraestrutura de saúde local limitada a atendimentos primários.
\item
  Dependência de centros maiores para insumos técnicos.
\end{itemize}

\subsubsection{Recomendação
Técnica}

Altamente indicado para estabelecimentos de autossuficiência tecnificada
em solo de alta performance. Requer investimento em segurança
patrimonial e autonomia em saúde. O GSS de 6.8 reflete a riqueza dos
recursos produtivos em equilíbrio com os desafios de segurança regional.
Referencia atual de score: GSS 6.8.

\subsubsection{}

\begin{itemize}
\tightlist
\item
  \begin{enumerate}
  \def\labelenumi{\arabic{enumi})}
  \tightlist
  \item
    Dados censitários validados (INE 2022).
  \end{enumerate}
\item
  \begin{enumerate}
  \def\labelenumi{\arabic{enumi})}
  \setcounter{enumi}{1}
  \tightlist
  \item
    Relatórios de segurança regional da SENAD e FTC.
  \end{enumerate}
\item
  \begin{enumerate}
  \def\labelenumi{\arabic{enumi})}
  \setcounter{enumi}{2}
  \tightlist
  \item
    Mapa de solos e aptidão agrícola de Amambay (MAG).
  \end{enumerate}
\item
  \begin{enumerate}
  \def\labelenumi{\arabic{enumi})}
  \setcounter{enumi}{3}
  \tightlist
  \item
    Registro de expansão infraestrutural do MOPC.
  \end{enumerate}
\end{itemize}}{dist:13_Amambay:Karapai}
\begin{table}[ht!]
\small \begin{tabular}{p{3.0cm}p{0.8cm}p{6.0cm}} \toprule
\textbf{Critério} & \textbf{Nota} & \textbf{Justificativa} \\ \midrule
A - Ameaças Estratégicas & 6.5 & Localização remota; boa invisibilidade tática, equilibrada pelo risco regional de segurança \\
B - Risco Social & 8.0 & Densidade populacional extremamente baixa; isolamento social superior \\
C - Riscos Naturais & 7.0 & Geologia muito estável e baixo risco de desastres naturais graves \\
D - Autossuficiência & 8.0 & Recursos agrícolas de elite e abundância hídrica subterrânea \\
E - Institucional & 4.5 & Suporte institucional em formação; serviços básicos funcionais mas limitados \\
\midrule \textbf{GSS FINAL} & \textbf{6.8} & \textbf{Moderadamente Seguro (Indicado para projetos de autossuficiência em solo de elite, exigindo preparo tático para o contexto regional).} \\ \bottomrule \end{tabular}
\end{table}
\subsubsection{Vulnerabilidades e
Mitigação}\label{vuln-Karapai-vulnerabilidades-e-mitigauxe7uxe3o}

\begin{itemize}
\tightlist
\item
  \textbf{Vulnerabilidade Tática Regional:} Insegurança ligada ao crime
  organizado no departamento (Mitigação: manutenção de sistemas de
  segurança patrimonial robustos e integração com redes de produtores
  locais).
\item
  \textbf{Dependência de Serviços Externos:} Saúde e logística
  especializada em cidades distantes (Mitigação: manutenção de estoque
  médico básico e kit trauma local).
\item
  \textbf{Fragilidade Elétrica:} Rede rural vulnerável (Mitigação:
  instalação de sistemas solares fotovoltaicos e geradores redundantes).
\end{itemize}

Notas (0-10): - A: 6.5 (Localização remota; boa invisibilidade tática,
equilibrada pelo risco regional de segurança) - B: 8.0 (Densidade
populacional extremamente baixa; isolamento social superior) - C: 7.0
(Geologia muito estável e baixo risco de desastres naturais graves) - D:
8.0 (Recursos agrícolas de elite e abundância hídrica subterrânea) - E:
4.5 (Suporte institucional em formação; serviços básicos funcionais mas
limitados)

Calculo: - GSS = ((A * 2.5) + (B * 2.0) + (C * 1.5) + (D * 2.0) + (E *
2.0)) / 10 - GSS: 6.8

Classificacao: - Moderadamente Seguro (Indicado para projetos de
autossuficiência em solo de elite, exigindo preparo tático para o
contexto regional).
\subsubsection{Dossiê de Campo}\label{dossie-Karapai-dossiuxea-de-campo}
\begin{description}[style=nextline,leftmargin=0.5cm,font=\bfseries]
\item[Coordenadas]  23°45'S 56°15'W.
\item[Topografia]  Relevo ondulado com mesetas mecanizáveis e remanescentes de floresta subtropical. Solo de altíssima produtividade (Terra Roxa). Altitude \textasciitilde 250m. Área de \textasciitilde 1.274 km². Geologicamente estável, risco sísmico desprezível.
\item[Alvos Estratégicos]  Áreas de produção agrícola intensiva (soja, milho); Eixo logístico da Ruta PY11 (Conexão Amambay-San Pedro); Reservas florestais. Ausência de alvos militares diretos, oferecendo boa invisibilidade estratégica, embora situada em departamento de vigilância tática permanente.
\item[Fallout]  Ventos predominantes do Norte (NE) e Leste. Baixo risco de fallout direto de alvos continentais primários.
\item[População]  5.382 habitantes (Censo 2022 Final). Município jovem, em fase de expansão demográfica ligada ao agronegócio.
\item[Densidade]  \textasciitilde 4,2 hab/km² (Densidade extremamente baixa, ideal para isolamento social tático e privacidade absoluta fora do núcleo urbano).
\item[Vias de Acesso]  Localização estratégica sobre a Ruta PY11. Conectividade asfáltica eficiente para o Sul (Santa Rosa del Aguaray) e Leste (Capitán Bado). Malha de estradas internas de terra que exigem veículos robustos em períodos chuvosos.
\item[Serviços]  Infraestrutura institucional básica. Unidade de Saúde da Família (USF) local. Dependência de Capitán Bado (50 km) ou Pedro Juan Caballero para serviços hospitalares de alta complexidade.
\item[Custo de Vida]  Baixo; economia baseada na exportação de grãos e pecuária.
\item[Preço da Terra]  US\$ 6.000-10.000/ha (terras mecanizadas de elite); US\$ 2.000-4.000/ha (áreas de pecuária/reserva).
\item[Hidrologia]  Baixo risco de inundações sistêmicas devido ao relevo ondulado e boa drenagem natural. Diversos arroios perenes de alta qualidade hídrica cruzam o distrito.
\item[Clima]  Tropical Úmido. Verões intensos e chuvosos (média 1.700 mm/ano); invernos amenos com geadas ocasionais.
\item[Energia]  Rede nacional (Itaipu); instável em áreas rurais remotas durante picos de demanda agrícola.
\item[Água]  Abundante via nascentes naturais e poços artesianos profundos; excelente qualidade hídrica subterrânea.
\item[Qualidade do Solo]  Latossolo Vermelho de Elite (Terra Roxa); solo profundo, estruturado e de altíssima fertilidade natural. Excelente aptidão para soja, milho e yerba mate.
\item[Recursos Locais]  Grande produção de grãos e gado de corte. Elevadíssimo potencial para autossuficiência alimentar básica em nível de propriedade devido à vasta área produtiva e baixa população.
\item[Segurança]  Zona rural pacífica em termos de crime comum urbano. No entanto, o departamento de Amambay exige vigilância tática elevada devido à presença de narcotráfico e grupos insurgentes em zonas de mata. Coesão social baseada na cultura produtora e comunidades indígenas locais.
\item[Leis Local: Município de criação recente (2013). Livre de Restrição de Fronteira]  Localizado fora da faixa de 50km da fronteira seca internacional.
\end{description}
\textbf{Fonte climática:} NASA POWER Climatology API (período 2001-2020)

\textbf{Fonte luminosa:} estimativa world atlas

\paragraph{Irradiação Solar
(kWh/m²/dia)}\label{irradiauxe7uxe3o-solar-kwhmuxb2dia}

{\def\LTcaptype{none} % do not increment counter
\begin{longtable}[]{@{}
  >{\raggedright\arraybackslash}p{(\linewidth - 24\tabcolsep) * \real{0.0746}}
  >{\raggedright\arraybackslash}p{(\linewidth - 24\tabcolsep) * \real{0.0746}}
  >{\raggedright\arraybackslash}p{(\linewidth - 24\tabcolsep) * \real{0.0746}}
  >{\raggedright\arraybackslash}p{(\linewidth - 24\tabcolsep) * \real{0.0746}}
  >{\raggedright\arraybackslash}p{(\linewidth - 24\tabcolsep) * \real{0.0746}}
  >{\raggedright\arraybackslash}p{(\linewidth - 24\tabcolsep) * \real{0.0746}}
  >{\raggedright\arraybackslash}p{(\linewidth - 24\tabcolsep) * \real{0.0746}}
  >{\raggedright\arraybackslash}p{(\linewidth - 24\tabcolsep) * \real{0.0746}}
  >{\raggedright\arraybackslash}p{(\linewidth - 24\tabcolsep) * \real{0.0746}}
  >{\raggedright\arraybackslash}p{(\linewidth - 24\tabcolsep) * \real{0.0746}}
  >{\raggedright\arraybackslash}p{(\linewidth - 24\tabcolsep) * \real{0.0746}}
  >{\raggedright\arraybackslash}p{(\linewidth - 24\tabcolsep) * \real{0.0746}}
  >{\raggedright\arraybackslash}p{(\linewidth - 24\tabcolsep) * \real{0.1045}}@{}}
\toprule\noalign{}
\begin{minipage}[b]{\linewidth}\raggedright
Jan
\end{minipage} & \begin{minipage}[b]{\linewidth}\raggedright
Fev
\end{minipage} & \begin{minipage}[b]{\linewidth}\raggedright
Mar
\end{minipage} & \begin{minipage}[b]{\linewidth}\raggedright
Abr
\end{minipage} & \begin{minipage}[b]{\linewidth}\raggedright
Mai
\end{minipage} & \begin{minipage}[b]{\linewidth}\raggedright
Jun
\end{minipage} & \begin{minipage}[b]{\linewidth}\raggedright
Jul
\end{minipage} & \begin{minipage}[b]{\linewidth}\raggedright
Ago
\end{minipage} & \begin{minipage}[b]{\linewidth}\raggedright
Set
\end{minipage} & \begin{minipage}[b]{\linewidth}\raggedright
Out
\end{minipage} & \begin{minipage}[b]{\linewidth}\raggedright
Nov
\end{minipage} & \begin{minipage}[b]{\linewidth}\raggedright
Dez
\end{minipage} & \begin{minipage}[b]{\linewidth}\raggedright
Média
\end{minipage} \\
\midrule\noalign{}
\endhead
\bottomrule\noalign{}
\endlastfoot
6.60 & 6.02 & 5.62 & 4.65 & 3.48 & 3.07 & 3.35 & 4.15 & 4.72 & 5.46 &
6.29 & 6.53 & \textbf{4.99} \\
\end{longtable}
}

\textbf{Inclinação solar recomendada:} 23° N (anual) · 33° N (inverno
jun-ago) · 13° N (verão nov-jan)

\paragraph{Precipitação (mm/mês)}\label{precipitauxe7uxe3o-mmmuxeas}

{\def\LTcaptype{none} % do not increment counter
\begin{longtable}[]{@{}
  >{\raggedright\arraybackslash}p{(\linewidth - 24\tabcolsep) * \real{0.0704}}
  >{\raggedright\arraybackslash}p{(\linewidth - 24\tabcolsep) * \real{0.0704}}
  >{\raggedright\arraybackslash}p{(\linewidth - 24\tabcolsep) * \real{0.0704}}
  >{\raggedright\arraybackslash}p{(\linewidth - 24\tabcolsep) * \real{0.0704}}
  >{\raggedright\arraybackslash}p{(\linewidth - 24\tabcolsep) * \real{0.0704}}
  >{\raggedright\arraybackslash}p{(\linewidth - 24\tabcolsep) * \real{0.0704}}
  >{\raggedright\arraybackslash}p{(\linewidth - 24\tabcolsep) * \real{0.0704}}
  >{\raggedright\arraybackslash}p{(\linewidth - 24\tabcolsep) * \real{0.0704}}
  >{\raggedright\arraybackslash}p{(\linewidth - 24\tabcolsep) * \real{0.0704}}
  >{\raggedright\arraybackslash}p{(\linewidth - 24\tabcolsep) * \real{0.0704}}
  >{\raggedright\arraybackslash}p{(\linewidth - 24\tabcolsep) * \real{0.0704}}
  >{\raggedright\arraybackslash}p{(\linewidth - 24\tabcolsep) * \real{0.0704}}
  >{\raggedright\arraybackslash}p{(\linewidth - 24\tabcolsep) * \real{0.1549}}@{}}
\toprule\noalign{}
\begin{minipage}[b]{\linewidth}\raggedright
Jan
\end{minipage} & \begin{minipage}[b]{\linewidth}\raggedright
Fev
\end{minipage} & \begin{minipage}[b]{\linewidth}\raggedright
Mar
\end{minipage} & \begin{minipage}[b]{\linewidth}\raggedright
Abr
\end{minipage} & \begin{minipage}[b]{\linewidth}\raggedright
Mai
\end{minipage} & \begin{minipage}[b]{\linewidth}\raggedright
Jun
\end{minipage} & \begin{minipage}[b]{\linewidth}\raggedright
Jul
\end{minipage} & \begin{minipage}[b]{\linewidth}\raggedright
Ago
\end{minipage} & \begin{minipage}[b]{\linewidth}\raggedright
Set
\end{minipage} & \begin{minipage}[b]{\linewidth}\raggedright
Out
\end{minipage} & \begin{minipage}[b]{\linewidth}\raggedright
Nov
\end{minipage} & \begin{minipage}[b]{\linewidth}\raggedright
Dez
\end{minipage} & \begin{minipage}[b]{\linewidth}\raggedright
Total/ano
\end{minipage} \\
\midrule\noalign{}
\endhead
\bottomrule\noalign{}
\endlastfoot
156 & 171 & 120 & 138 & 129 & 81 & 54 & 43 & 85 & 166 & 195 & 170 &
\textbf{1513 mm} \\
\end{longtable}
}

\paragraph{Poluição Luminosa}\label{poluiuxe7uxe3o-luminosa}

{\def\LTcaptype{none} % do not increment counter
\begin{longtable}[]{@{}ll@{}}
\toprule\noalign{}
Parâmetro & Valor \\
\midrule\noalign{}
\endhead
\bottomrule\noalign{}
\endlastfoot
Escala Bortle & 4 --- Céu rural-suburbano \\
Radiância artificial & 2.0 nW/cm²/sr \\
\end{longtable}
}
\paragraph{Solo (SoilGrids 2.0, média ponderada 0--30
cm)}

{\def\LTcaptype{none} % do not increment counter
\begin{longtable}[]{@{}ll@{}}
\toprule\noalign{}
Parâmetro & Valor \\
\midrule\noalign{}
\endhead
\bottomrule\noalign{}
\endlastfoot
pH (H\(_2\)O) & 5.3 \\
Carbono orgânico (SOC) & 17.51 g/kg \\
Argila & 27.46 \% \\
Areia & 52.04 \% \\
Silte (calc.) & 20.5 \% \\
Densidade aparente & 1.34 g/cm³ \\
\textbf{Aptidão agrícola} & \textbf{Média} \\
\end{longtable}
}

Fonte: ISRIC SoilGrids 2.0 via WCS. Coords: -23.75°, -56.25°. Média
ponderada camadas 0-5, 5-15, 15-30 cm.

\begin{itemize}
\tightlist
\item
  \textbf{Homicidios/100k:} \textasciitilde25,0 (Regional/Amambay -
  Concentrado em áreas de fronteira).
\end{itemize}
\subsubsection{Indicadores Sociais}

\textbf{Fontes:} PNUD 2020, INE Censo 2022, INE EPHC 2023, Ministerio
Público 2024\\
\textbf{Nota:} Valores marcados como \emph{(dept.)} referem-se ao
departamento de Amambay; valores \emph{(dist.)} são específicos deste
distrito. Dados marcados \emph{(est.)} são estimativas.

{\def\LTcaptype{none} % do not increment counter
\begin{longtable}[]{@{}
  >{\raggedright\arraybackslash}p{(\linewidth - 4\tabcolsep) * \real{0.4231}}
  >{\raggedright\arraybackslash}p{(\linewidth - 4\tabcolsep) * \real{0.2692}}
  >{\raggedright\arraybackslash}p{(\linewidth - 4\tabcolsep) * \real{0.3077}}@{}}
\toprule\noalign{}
\begin{minipage}[b]{\linewidth}\raggedright
Indicador
\end{minipage} & \begin{minipage}[b]{\linewidth}\raggedright
Valor
\end{minipage} & \begin{minipage}[b]{\linewidth}\raggedright
Âmbito
\end{minipage} \\
\midrule\noalign{}
\endhead
\bottomrule\noalign{}
\endlastfoot
IDH (2020) & 0,720 & dept. \\
Ranking IDH nacional & 5/18 & dept. \\
Esperança de vida & 73,6 anos & dept. \\
Escolaridade média & 6.9 anos & dist. \\
RNB per capita (USD PPA) & 9.500 & dept. \\
Pobreza monetária (\%) & 22,8\% & dept. \\
Pobreza extrema (\%) & 4,5\% & dept. \\
Índice de Gini & 0,448 & dept. \\
Acesso a água potável (\%) & 82,1\% & dept. \\
Acesso a saneamento (\%) & 80,5\% & dept. \\
Taxa de homicídios (est., /100k hab) & \textasciitilde30--50 (estimativa
2022--2024; histórico: 107,4 em 2016, 69,9 em 2017) & dept. \\
Índice de segurança & baixo & dept. \\
População (Censo 2022) & 5.382 habitantes & dist. \\
Idade mediana & N/D & dist. \\
Taxa de fecundidade & N/D & dist. \\
\end{longtable}
}
\subsubsection{Combustível}

\textbf{Referência:} PETROPAR / postos locais (2024) \textbf{Tipo de
localidade:} Interior

{\def\LTcaptype{none} % do not increment counter
\begin{longtable}[]{@{}lll@{}}
\toprule\noalign{}
Combustível & USD/litro & Gs/litro (aprox.) \\
\midrule\noalign{}
\endhead
\bottomrule\noalign{}
\endlastfoot
Gasolina 93 oct & 1.00 & 7,400 \\
Gasolina 97 oct (premium) & 1.10 & 8,140 \\
Diesel & 0.93 & 6,882 \\
\end{longtable}
}

\begin{quote}
Preços podem variar ±5\% conforme posto e sazonalidade. Chaco e interior
remoto apresentam maior variação.
\end{quote}

\subsubsection{Cobertura Celular}

\textbf{Fonte:} CONATEL PY / operadoras (2024)

{\def\LTcaptype{none} % do not increment counter
\begin{longtable}[]{@{}ll@{}}
\toprule\noalign{}
Parâmetro & Valor \\
\midrule\noalign{}
\endhead
\bottomrule\noalign{}
\endlastfoot
Cobertura 4G população (dept.) & 85\% \\
Cobertura 4G área rural & 68\% \\
Melhor operadora & Tigo \\
Qualidade rural & moderada \\
\end{longtable}
}

\begin{quote}
Para áreas rurais fora do núcleo urbano, recomenda-se chip Tigo como
principal e Personal como backup.
\end{quote}

\subsubsection{Internet}

\textbf{Fonte:} CONATEL / Speedtest Ookla (2024)

{\def\LTcaptype{none} % do not increment counter
\begin{longtable}[]{@{}ll@{}}
\toprule\noalign{}
Parâmetro & Valor \\
\midrule\noalign{}
\endhead
\bottomrule\noalign{}
\endlastfoot
Velocidade média download & 55 Mbps \\
Domicílios com internet (dept.) & 62\% \\
Tecnologia predominante & rádio/fibra \\
Opção rural & Starlink disponível (\textasciitilde USD 44/mês) \\
\end{longtable}
}

\subsubsection{Mercado Imobiliário e Terra
Rural}

\textbf{Fonte:} INDERT / Clasificados.com.py (2024)

{\def\LTcaptype{none} % do not increment counter
\begin{longtable}[]{@{}ll@{}}
\toprule\noalign{}
Tipo & Referência \\
\midrule\noalign{}
\endhead
\bottomrule\noalign{}
\endlastfoot
Terra agrícola alta prod. (USD/ha) & 5,000 \\
Imóvel urbano (USD/m²) & 900 \\
Aluguel 2 quartos (USD/mês) & 380 \\
\end{longtable}
}

\begin{quote}
Valores de referência departamental. Localidades menores podem ter
preços 20--40\% abaixo da capital departamental.
\end{quote}

\subsubsection{Saúde}

\textbf{Fonte:} MSPBS / IPS Paraguay (2024-2026), consolidação
departamental e proxy local

{\def\LTcaptype{none} % do not increment counter
\begin{longtable}[]{@{}ll@{}}
\toprule\noalign{}
Serviço & Disponibilidade \\
\midrule\noalign{}
\endhead
\bottomrule\noalign{}
\endlastfoot
USF / Posto de Saúde & sim \\
Hospital Regional & sim \\
IPS (seguro social) & não \\
Farmácia & sim \\
Distância ao hospital de referência & \textasciitilde134 km (Pedro Juan
Caballero) \\
\end{longtable}
}

\textbf{Principais estabelecimentos:} USF local; referencia hospitalar
em Pedro Juan Caballero

\textbf{Observação para imigrantes:} Atendimento primario local ou em
raio curto; casos de maior complexidade seguem para Pedro Juan
Caballero. Cobertura privada continua recomendada para especialidades e
urgências de maior complexidade.
\newpage
